% TOP_PROBLEM: {"id": 3252, "title": "List Colourings of Complete Multipartite Graphs with 2 Big Parts", "category_id": 3, "category": {"id": 3, "name": "graph_theory", "display_name": "Graph Theory", "description": "Problems involving graphs, networks, and their properties.", "slug": "graph-theory", "order_index": 3, "created_at": "2026-07-31T15:26:25.671Z"}, "status": "open", "problem_number": "OPG-59927", "set_id": 12}
% FIRST_POSTED: 2026-10-01
% ENTRY_KIND: statement_only
% UnsolvedMath ID 3252; source catalogue code OPG-59927
% !TeX program = lualatex
% Definition and sources only; this file is not an LLM solution attempt.
\documentclass[11pt,a4paper]{article}
\usepackage[margin=25mm]{geometry}
\usepackage{fontspec}
\setmainfont{lmroman10-regular.otf}[BoldFont=lmroman10-bold.otf,
 ItalicFont=lmroman10-italic.otf,BoldItalicFont=lmroman10-bolditalic.otf]
\usepackage{amsmath,amssymb,mathtools}
\usepackage{unicode-math}
\setmathfont{latinmodern-math.otf}
\AtBeginDocument{\let\setminus\smallsetminus}
\usepackage{xurl,hyperref}
\hypersetup{colorlinks=true,urlcolor=blue}
\setcounter{secnumdepth}{0}
\setlength{\parindent}{0pt}
\setlength{\parskip}{5pt}
\setlength{\emergencystretch}{3em}
\begin{document}
\section*{List Colourings of Complete Multipartite Graphs with 2 Big Parts}\label{3252}
{\small UnsolvedMath ID 3252 · OPG-59927 · Graph Theory · OpenGarden}
\subsection{Definitions and mathematical statement}
Fix integers $a,b\ge2$. The complete bipartite graph $K_{a,b}$ has disjoint
independent vertex sets of sizes $a$ and $b$, with every edge between the
two sets. The join $G+H$ of disjoint graphs adds every edge between
$V(G)$ and $V(H)$. Write $K_t$ for the complete graph on $t$ vertices,
including $K_0$ with no vertices. Thus $K_{a,b}+K_t$ is a complete
multipartite graph with part sizes $a,b,1,\ldots,1$.

A proper colouring assigns different colours to adjacent vertices.
The chromatic number of this join is $t+2$: the two large parts require
two colours, and each of the $t$ joined clique vertices requires its own
additional colour. Its list chromatic number $\chi_\ell$ is the least
integer $q$ such that every assignment of $q$ permitted colours to each
vertex admits a proper colouring from those lists.

Determine, for every $a,b\ge2$, the parameter
\[
 \phi(a,b)=\min\{t\in\mathbb N_0:
                 \chi_\ell(K_{a,b}+K_t)=t+2\}.
\]
The source discusses a theorem ensuring the defining set is nonempty.
A complete determination must supply both a universal list-colouring
argument at the proposed $t$ and, for every smaller $t$, a list assignment
that cannot be coloured with $t+2$ available colours per vertex. Ordinary
chromatic number alone does not decide the threshold. Interchanging $a$
and $b$ leaves the graph unchanged, so the sought function is symmetric.
\subsection{Short English statement}
How many mutually adjacent universal vertices must be added to a complete bipartite graph before its choice number equals its chromatic number?
\subsection{Sources}
\begin{itemize}
\item[S1] ULAM.AI and the UnsolvedMath Contributors, supplied problems.json, record 3252 (OPG-59927), OpenGarden collection. Catalogue identity and source transcription; CC BY 4.0.
\url{https://huggingface.co/datasets/ulamai/UnsolvedMath}.
\item[S2] Open Problem Garden, List Colourings of Complete Multipartite Graphs with 2 Big Parts. Original problem and discussion; the mathematical formulation below is expanded from the supplied source record.
\url{http://www.openproblemgarden.org/op/list_colourings_of_complete_multipartite_graphs_with_2_big_parts}.
\end{itemize}
\end{document}
